% CP-VI-0033
% target_chapter: VI/27 — Integral Schemes and Function Fields
% source_unit_id: IMP-DL-D27460A55B-P002
% source: Downloads/theory-of-algebraic-geometry-5.tex L278-L462
% solution_provenance: SOURCE_EXPLICIT_SOLUTION

\subsection*{Problem 3 - Integral Schemes, Reduced Schemes, and Irreducible Schemes}

\textbf{Problem.}
We say a scheme \(X\) is \emph{irreducible} if it is irreducible as a topological space, i.e. whenever
\[
X=X_1\cup X_2
\]
with \(X_1,X_2\) closed subsets, then either \(X_1=X\) or \(X_2=X\).

We say a scheme \(X\) is \emph{reduced} if for every open subset \(U\subseteq X\), the ring
\[
\mathcal O_X(U)
\]
has no nilpotents.

We say a scheme \(X\) is \emph{integral} if for every open subset \(U\subseteq X\), the ring
\[
\mathcal O_X(U)
\]
is an integral domain.

Show that a scheme is integral if and only if it is reduced and irreducible.

\textbf{Solution.}

We prove both implications.

\subsection*{Integral implies reduced and irreducible}

Assume that \(X\) is integral. By definition, for every open subset \(U\subseteq X\), the ring
\[
\mathcal O_X(U)
\]
is an integral domain. Every integral domain is reduced. Therefore \(\mathcal O_X(U)\) has no nilpotents for every open \(U\), so \(X\) is reduced.

It remains to prove that \(X\) is irreducible.

Recall that a topological space is irreducible if and only if every two nonempty open subsets intersect. Suppose, for contradiction, that \(X\) is not irreducible. Then there exist nonempty open subsets
\[
U_1,U_2\subseteq X
\]
such that
\[
U_1\cap U_2=\varnothing.
\]
Consider the open subset
\[
W=U_1\cup U_2.
\]
Since \(U_1\) and \(U_2\) are disjoint open subsets, the sheaf property gives
\[
\mathcal O_X(W)
\cong
\mathcal O_X(U_1)\times \mathcal O_X(U_2).
\]
But the product of two nonzero rings has zero divisors. Indeed,
\[
(1,0)(0,1)=(0,0),
\]
while
\[
(1,0)\neq 0,
\qquad
(0,1)\neq 0.
\]
Therefore \(\mathcal O_X(W)\) is not an integral domain. This contradicts the assumption that \(X\) is integral.

Hence every two nonempty open subsets of \(X\) intersect, so \(X\) is irreducible.

Therefore
\[
X\text{ integral}
\implies
X\text{ reduced and irreducible}.
\]

\subsection*{Reduced and irreducible implies integral}

Now assume that \(X\) is reduced and irreducible. We must show that for every open subset \(U\subseteq X\),
\[
\mathcal O_X(U)
\]
is an integral domain.

First observe that every nonempty open subset of an irreducible topological space is irreducible. Hence every nonempty open subset \(U\subseteq X\) is irreducible.

Let \(U\subseteq X\) be a nonempty open subset. Take
\[
s,t\in \mathcal O_X(U)
\]
such that
\[
st=0.
\]
We want to show that either \(s=0\) or \(t=0\).

Suppose, for contradiction, that
\[
s\neq 0
\qquad\text{and}\qquad
t\neq 0.
\]
Then there exist points \(x,y\in U\) such that
\[
s_x\neq 0,
\qquad
t_y\neq 0.
\]
The nonvanishing locus of a section is open, so the sets
\[
U_s=\{p\in U : s_p\neq 0\},
\qquad
U_t=\{p\in U : t_p\neq 0\}
\]
are nonempty open subsets of \(U\).

Since \(U\) is irreducible, two nonempty open subsets of \(U\) must intersect. Therefore
\[
U_s\cap U_t\neq \varnothing.
\]
Choose
\[
p\in U_s\cap U_t.
\]
Then
\[
s_p\neq 0,
\qquad
t_p\neq 0.
\]

Now choose an affine open neighborhood
\[
W=\operatorname{Spec} R
\]
of \(p\) contained in \(U\). Since \(X\) is reduced and irreducible, the open subscheme \(W\) is also reduced and irreducible. Therefore \(R\) is a reduced ring and \(\operatorname{Spec} R\) is irreducible.

For an affine scheme \(\operatorname{Spec} R\), irreducibility of \(\operatorname{Spec} R\) means that the nilradical of \(R\) is a prime ideal. Since \(R\) is reduced, its nilradical is zero. Hence
\[
(0)\subseteq R
\]
is prime, so \(R\) is an integral domain.

Thus the local ring
\[
\mathcal O_{X,p}
\]
is a localization of an integral domain, hence is itself an integral domain.

But
\[
(st)_p=s_p t_p=0
\]
in \(\mathcal O_{X,p}\). Since \(\mathcal O_{X,p}\) is an integral domain, this implies
\[
s_p=0
\qquad\text{or}\qquad
t_p=0,
\]
contradicting the choice of \(p\).

Therefore our assumption was false, and so either
\[
s=0
\qquad\text{or}\qquad
t=0.
\]
Hence \(\mathcal O_X(U)\) has no zero divisors.

Since \(X\) is reduced, \(\mathcal O_X(U)\) also has no nilpotents. Therefore \(\mathcal O_X(U)\) is an integral domain.

Thus \(X\) is integral.

\subsection*{Conclusion}

We have proved both directions, so
\[
\boxed{
	X\text{ is integral}
	\iff
	X\text{ is reduced and irreducible}.
}
\]
